%document class
\documentclass[10pt,oneside]{book}
%%%% Page Info + Commands %%%%%{

%packages
\usepackage{pgfplots}
\pgfplotsset{compat=1.18}
\usepackage{amsmath}
\usepackage{amsfonts}
\usepackage{charter}
\usepackage{amsthm,letltxmacro}
\usepackage{amssymb}
\usepackage[most]{tcolorbox}
\usepackage{enumerate}
\usepackage{enumitem}
\usepackage[dvipsnames]{xcolor}
\usepackage{changepage}
\usepackage{mdframed}

%This will shrink the page
\newcommand{\smallpage}{  \setlength \oddsidemargin{.25in}
            \setlength \textwidth{5.5in}
            \setlength \topmargin{0in}
            \setlength \textheight{9in}}


% Commands for sets of numbers
\newcommand{\Z}{\mathbb{Z}}\newcommand{\R}{\mathbb{R}}\newcommand{\N}{\mathbb{N}}

% gordie spacing and boxing commands
\newcommand{\gspc}{\vspace{0.13cm}} % 0.5 space
\newcommand{\gline}{\gspc\\} % 1.5 space
\newcommand{\gbline}{\gspc\gspc\\} % double space
\newcommand{\gbox}[1]{\fbox{\parbox{\linewidth}{#1}}} % multiline box command
\newcommand{\gmod}[1]{\,(\text{mod}\,#1)}


\newcommand{\gimage}[3]{
\begin{figure}[h]   
    \centering          
    \includegraphics[width=#2\textwidth]{#1}  
    \caption{#3}
    \label{fig:#1}
\end{figure}\\
}

\newcommand{\centerit}[1]{{\begin{center} #1 \end{center}}}
\newcommand{\ul}[1]{\underline{#1}}

%%%%%%%% command for graphics %%%%%%%%%%%%%
\usepackage{fancyhdr}
%}

%%%% Page Setup %%%%%%%
\smallpage\sloppy
\pagestyle{fancy}
\lhead{\large{\textbf{Jan 23rd - Notes}}}
%\rfoot{\thepage/\pageref{LastPage} }
\setlength{\headheight}{10pt}

%coloring with color
\newmdenv[
  backgroundcolor=gray!8,
  linewidth=0pt,
  innerleftmargin=0pt,
  innerrightmargin=0pt,
  skipabove=\baselineskip,
  skipbelow=\baselineskip
]{coloredadjust}

% Proof writing
\LetLtxMacro\oldproof\proof
\let\endoldproof\endproof
\renewenvironment{proof}[1][\proofname]
  {\vspace{0.2cm}\begin{adjustwidth}{2em}{2em}
   \oldproof[#1]}
  {\endoldproof
\end{adjustwidth}}


% problem writing
\newenvironment{problem}[1]
  {\vspace{-0.1cm}\begin{coloredadjust}\begin{adjustwidth}{2em}{2em}
   \textit{Problem. }#1}
  {\endoldproof
\end{adjustwidth}\end{coloredadjust}\gspc}

% theorem writing
\newtcbtheorem[
]{theorem}{Theorem}%
{%
  enhanced,
  colback=blue!3,
  colframe=blue!60!black,
  coltitle=blue!90!black,
  fonttitle=\bfseries,
  boxrule=0pt,
  toprule=2pt,
  bottomrule=0pt,
  leftrule=0pt,
  rightrule=0pt,
  arc=3pt,
  outer arc=3pt,
  attach title to upper,
  title={#1},
}{thm}

\newtcbtheorem[
]{criterion}{Criterion}%
{%
  enhanced,
  colback=green!3,
  colframe=green!60!black,
  coltitle=green!40!black,
  fonttitle=\bfseries,
  boxrule=4pt,
  toprule=2pt,
  bottomrule=1pt,
  leftrule=0pt,
  rightrule=0pt,
  arc=3pt,
  outer arc=3pt,
  attach title to upper,
  title={#1},
}{ctn}

\newtcbtheorem[
]{corollary}{Corollary}%
{%
  enhanced,
  colback=yellow!2,
  colframe=yellow!60!black,
  coltitle=yellow!40!black,
  fonttitle=\bfseries,
  boxrule=4pt,
  toprule=2pt,
  bottomrule=1pt,
  leftrule=0pt,
  rightrule=0pt,
  arc=3pt,
  outer arc=3pt,
  attach title to upper,
  title={#1},
}{ctn}


\begin{document}

\newcommand{\cg}[1]{\color{OliveGreen}#1\color{white}}
\newcommand{\cb}[1]{\color{BlueViolet}#1\color{white}}

\subsection*{Remember from Last Time}
\subsubsection*{Euler's Criterion}
\begin{criterion}{Euler's Criterion}{}\gline
  If $p$ is an odd prime, and $p\nmid a$:
  \begin{itemize}[itemsep=0pt, parsep=0pt]
    \item $a$ is a quadratic residue of $p \Leftrightarrow a^{(p-1)/2}\equiv 1\gmod{p}$.
    \item $a$ is a quadratic nonresidue of $p\Leftrightarrow a^{(p-1)/2}\equiv -1 \gmod p$
  \end{itemize}
\end{criterion}
We rephrase using the Legendre symbols: \gline
$(a/p)\equiv a^{(p-1)/2}\gmod p$\\
Multiplicativity: $(ab/p) = (a/p)(b/p)$. (i.e. $ab$ is a quadratic residue if $!(a \texttt{ XOR } b))$.
\begin{corollary}{Euler's Criterion}{}\gline
  We also have the following properties:
  \begin{itemize}[itemsep=0pt, parsep=0pt]
    \item $(1/p) = 1$ for all $p$
    \item $(4/p) = 1$ for all $p$
    \item $(-1/p) = 1$ if and only if $p\equiv 1 \gmod 4$
  \end{itemize}
\end{corollary}

This brings us to Dirchlet's Theorem:
\begin{theorem}{Dirchlet's Theorem }{}\gline
  Let $\gcd(a,b) =1$. There are $\infty$-many primes $p$ for which $p\equiv b \gmod{a}$.\gline
  This leads to the corollary such that there are $\infty$ many primes such that $p\equiv 1 \gmod 4$. 
  \begin{proof}
    Assume there are finitely many primes. Create $X = 4p_1^2p_2^2\cdots p_r^2 + 1$.\gline
    Facts: $X\equiv 1 \gmod 4$, and $X$ has at lease one prime factor: $q$. 
    \begin{itemize}
      \item No $p_i\mid X$, otherwise $p_i\mid 1$, so that means that $q\not\equiv 1 \gmod 4$, so $q\equiv 3\gmod 4$. Now we manipulate with algebra:
      \begin{align*}
        X&\equiv 0 \gmod q\\
        4p_1^2p_2\cdots p_r^2 +1&\equiv 0 \gmod q\\
        (2p_1p_2\cdots p_r)^2 + 1&\equiv 0 \gmod q\\
        (2p_1p_2\cdots p_r)^2 &\equiv -1 \gmod q
      \end{align*}
      Thus, that means that $-1$ is a quadratic residue $\gmod q$, so $q\equiv 1 \gmod 4$ via the corollary to Euler's Criterion, a contradiction. 
    \end{itemize}
  \end{proof}
\end{theorem}
Take the following example:
\begin{problem}
  Is 120 a quadratic residue for 239?\gline
  To use Euler's criterion, you'd need to find $120^{119}\gmod{239}$, which sounds miserable.\gline
  By multiplicity, we have that: $(120/239)=(2/239)(2/239)(2/239)(3/239)(5/239)$. Because $(2/239)$ and $(3/239)$ are both 1, we just needto find $(5/239).$
\end{problem}
\begin{theorem}{Quadratic Reciprocity Law} {}\gline
  Let $p,q$ be distinct odd primes. Then:
  \begin{gather*}
    (p/q)(q/p)= (-1)^{\left(\frac{p-1}2\right)\left(\frac{(q-1)}2\right)}
  \end{gather*}
\end{theorem}
\begin{problem}
  Let $p = 3, q= 239$. Then, we have that:
  \begin{gather*}
    (3/239)(239/3) = (-1)^{\left(\frac{3-1}2\right)\left(\frac{239-1}2\right)}
  \end{gather*}
  Well, we know that $239 \equiv 2 \gmod 3$, which is not a quadratic residue, so we know that $(239/3)$ is -1, so that must imply that $(3/239)=1$. Isn't that neat?
\end{problem}
Let's try another problem now:
\begin{problem}
  For which primes is 3 a quadratic residue?\gline
  Via the Quadratic Reciprocity Law, wehave that:
  \begin{align*}
    (3/q)(q/3)&=(-1)^{\left(\frac{3-1}2\right)\left(\frac{q-1}2\right)}\\
    &=(-1)^{\left(\frac{q-1}2\right)}
  \end{align*}
\end{problem}

\end{document}
